\documentclass[12pt,aspectratio=43]{beamer}

%%%%%%%%%%%%%%%%%%%%%%%%%%%%%%%%%%%%%%%%%%%%%%%%%%
%   Paquetes de configuración basica del texto   %
%%%%%%%%%%%%%%%%%%%%%%%%%%%%%%%%%%%%%%%%%%%%%%%%%%
\usepackage[spanish,es-nodecimaldot]{babel}
\usepackage[T1]{fontenc}
\usepackage{listings}
\usepackage{hyperref}
\usepackage{hyperxmp}
\usepackage{ifthen}

%%%%%%%%%%%%%%%%%%%%%%%%%%%%%%%%%%%%%%%%%%%%%%%%%%
%                   [listings]                   %
%   Configuración de vizualizacion del código    %
%%%%%%%%%%%%%%%%%%%%%%%%%%%%%%%%%%%%%%%%%%%%%%%%%%
\lstset{
	backgroundcolor=\color{lightgray!10},
	basicstyle=\ttfamily,
	extendedchars=true,
	showspaces=false,
	showstringspaces=false,
	captionpos=b,
	keywordstyle=\bfseries\color{cyan},
	commentstyle=\color{gray},
	stringstyle=\color{orange},
	escapeinside={!>}{<!},
	language=[LaTeX]TeX }

%%%%%%%%%%%%%%%%%%%%%%%%%%%%%%%%%%%%%%%%%%%%%%%%%%
%     Paquetes de configuración graficos del     %
%                   documento                    %
%%%%%%%%%%%%%%%%%%%%%%%%%%%%%%%%%%%%%%%%%%%%%%%%%%
\usepackage{graphicx}
\usepackage{tikz}
\usetikzlibrary{babel,shadows}
\usepackage{xcolor}
\usepackage{pdfpages}
\usepackage{animate}

\usepackage{ifxetex}
\ifxetex
	\usepackage[no-math]{fontspec}
	\setmainfont{AncizarSans}[
		Path = ../Fuentes/,
		Extension = .otf,
		BoldFont = *-B,
		ItalicFont = *-I,
		BoldItalicFont = *-BI ]
	\setmonofont{UbuntuMono}[
		Path = ../Fuentes/,
		Extension = .ttf,
		BoldFont = *-B,
		ItalicFont = *-I,
		BoldItalicFont = *-BI ]
	\newcommand{\lmr}{\fontfamily{lmr}\selectfont}
	\newcommand{\lmss}{\fontfamily{lmss}\selectfont}
	\newcommand{\lmtt}{\fontfamily{lmtt}\selectfont}
\fi

%%%%%%%%%%%%%%%%%%%%%%%%%%%%%%%%%%%%%%%%%%%%%%%%%%
%           Configuraciones de Beamer            %
%%%%%%%%%%%%%%%%%%%%%%%%%%%%%%%%%%%%%%%%%%%%%%%%%%
\definecolor{Igreen}{RGB}{148,180,59}
\definecolor{Ired}{RGB}{166,24,49}

\setbeamertemplate{navigation symbols}{}

\setbeamertemplate{itemize items}{\color{Igreen}\raisebox{.45ex}{\rule{.6ex}{.6ex}}}

\usefonttheme{professionalfonts}
\usefonttheme{serif}
\setbeamercolor{frametitle}{fg=Ired,bg=Igreen!50}
\setbeamercolor{alerted text}{fg=Ired}
\setbeamertemplate{enumerate item}{\color{Igreen}\bfseries\arabic{enumi}.}
\setbeamertemplate{enumerate subitem}{\color{Igreen}\bfseries\arabic{enumi}.\arabic{enumii}.}

\makeatletter
\newcommand{\ifratio}[2]{
	\ifthenelse
	{\lengthtest{\beamer@paperwidth=16cm} \AND \lengthtest{\beamer@paperheight=9cm}}
	{#1}
	{#2} }
\makeatother

\setbeamertemplate{title page}{
	\begin{tikzpicture}[overlay,remember picture]
	\path (current page.north west) ++(0.5cm,-0.5cm) node[below right] {\includegraphics[height=1.25cm]{../Escudo_UN}};
	\path (current page.north east) ++(-0.5cm,-0.5cm) node[below left] {\parbox[c][1.25cm][c]{\widthof{\Large \insertsubtitle}}{\Large \insertsubtitle}};
	\end{tikzpicture}
	
	\vfill
	\begin{center}
		{\Huge\inserttitle}\\
		\bigskip
		\insertauthor
	\end{center}
	\vfill
}

%%%%%%%%%%%%%%%%%%%%%%%%%%%%%%%%%%%%%%%%%%%%%%%%%%
%         Información de la presentación         %
%%%%%%%%%%%%%%%%%%%%%%%%%%%%%%%%%%%%%%%%%%%%%%%%%%

\title{Introducción a {\lmr\LaTeX}}
\subtitle{Curso Libre de {\lmr\LaTeX}}
\author{Joar Esteban Buitrago Carrillo}
\institute{Universidad Nacional de Colombia}
\date{}
\logo{\includegraphics[height=0.75cm]{../Escudo_UN}}
\hypersetup{
	pdfcopyright={Esta obra está bajo una licencia de Creative Commons Reconocimiento-CompartirIgual 4.0 Internacional.},
	pdflicenseurl={https://creativecommons.org/licenses/by-sa/4.0/deed.es}
}

\begin{document}
\begin{frame}[plain]
\titlepage
\end{frame}

{
\setbeamercolor{background canvas}{bg=}
\begin{frame}
\ifratio
	{\includepdf[pages=2]{../Licencia_Diapositiva.pdf}}
	{\includepdf[pages=1]{../Licencia_Diapositiva.pdf}}
\end{frame}
}

\begin{frame}{¿Que es {\lmr\LaTeX}?}{}
{\lmr\LaTeX} es un sistema de composición de textos de calidad artística y tipográfica. Es usado muy habitualmente en el mundo científico debido a todas sus características y facilidad de trabajo.\pause\\[2em]

{\lmr\LaTeX} es un conjunto de macros de {\lmr\TeX}, licenciado bajo la licencia {\it{\lmr L}PPL}, {\it{\lmr\LaTeX} Project Public License}. Por lo cual es {\em Software Libre} y gratuito.
\end{frame}

\begin{frame}[fragile]{Funciones Seno y Coseno}
\begin{center}
%	\begin{animateinline}[autoplay,loop,nomouse,poster=first]{24}
%		\multiframe{359}{nAngle=0+1}{
%			\begin{tikzpicture}
%			\draw[thick,-latex,blue] (-3,0) -- (3,0) node[right] {$x$};
%			\draw[thick,-latex,blue] (0,-3) -- (0,3) node[above] {$y$};
%			\draw[red,thick] (0,0) circle (2.5);
%			\draw[ultra thick,cyan] (0,0) -- (0,0 |- \nAngle:2.5) node[midway,left] {$\sin(\alpha)$};
%			\draw[ultra thick,orange] (0,0) -- (\nAngle:2.5 |- 0,0) node[midway,below] {$\cos(\alpha)$};
%			\draw[densely dotted,orange] (\nAngle:2.5) -- (\nAngle:2.5 |- 0,0);
%			\draw[densely dotted,cyan] (\nAngle:2.5) -- (0,0 |- \nAngle:2.5);
%			\draw[thick,green] (0.5,0) arc (0:\nAngle:0.5);
%			\draw[thick,green,opacity=0] (0.7,0) arc (0:\nAngle:0.7) node[midway,opacity=1] {$\alpha$};
%			\draw[ultra thick,red,-latex] (0,0) -- (\nAngle:2.5);
%			\end{tikzpicture}
%		}
%	\end{animateinline}
\end{center}
\end{frame}

\section{Clases}
\begin{frame}{Clases del documento}{}
Al trabajar en un documento escrito con {\lmr\LaTeX} debemos tener en cuenta la clase del documento en el cual queremos trabajar, la clase más usada generalmente es \texttt{\textbf{article}}.\pause\\[2em]

Las clases «estándares» de {\lmr\LaTeX} son:
\begin{itemize}
	\item {\color{Ired}\texttt{article}}: para artículos.
	\item {\color{Ired}\texttt{report}}: para reportes.
	\item {\color{Ired}\texttt{book}}: para libros.
	\item {\color{Ired}\texttt{letter}}: para las cartas.
	\item {\alt<3>{\color{black!30!Igreen}\textbf{\texttt{beamer}}}{\color{Ired} \texttt{slides}}}: para las presentaciones.
\end{itemize}
\end{frame}

\begin{frame}[fragile]{Clases del documento}{}
Para indicar la clase del documento a {\lmr\LaTeX} solo debemos de escribir en nuestro código:
\begin{center}
	\lstinline|\documentclass{<clase>}|
\end{center}
\pause
Donde \texttt{<clase>} es cualquiera de las clases ya mencionadas, u otras personalizadas o descargadas.
\end{frame}

\section{Estructura del código}
\begin{frame}[fragile]{Estructura del código}{}
La estructura básica del código de cualquier documento es muy sencilla, siempre es la misma.\pause

\begin{enumerate}
	\item Primero se coloca el comando \lstinline|\documentclass{<clase>}| con la clase escogida.\pause
	\item Desde aquí encontramos el \alert{\bf preámbulo} hasta que se encuentra la orden \lstinline|\begin{document}|.\pause
	\item Después tenemos el \alert{\bf cuerpo del documento} hasta encontrar la orden \lstinline|\end{document}|, todo lo que se encuentra después de esto es ignorado por {\lmr\LaTeX}.\pause
\end{enumerate}

Se debe tener cuidado al escribir, ya que existen comandos los cuales solo funcionan en una sección del código.
\end{frame}

\begin{frame}[fragile]{Estructura del código}{Código}
De manera más visual esto seria:
\begin{lstlisting}
\documentclass{<clase>}

Preámbulo ...

\begin{document}

Documento ...

\end{document}
\end{lstlisting}
\end{frame}

\section{Títulos}
\begin{frame}[fragile]{Títulos}{}
Al empezar a escribir, con cualquiera de las clases, necesitamos colocar el titulo de nuestro documento.\pause\\[2em]

Esto lo logramos mediante las ordenes:
\begin{itemize}
	\item \lstinline|\title{<título>}|
	\item \lstinline|\author{<autor o autores>}|
	\item \lstinline|\date{<fecha>}|
\end{itemize}\pause\vspace{2em}

Y para colocarlos en el documento de forma general, podemos hacerlo mediante la orden \lstinline|\maketitle|
\end{frame}

\begin{frame}[fragile]{Títulos}
Al escribir esto ya en un documento con la clase \alert{\bf articulo}, de manera normal, podemos obtener esto.\\[1em]

\begin{columns}
	\begin{column}{0.5\linewidth}
		{\small
\begin{lstlisting}
\documentclass{article}

\title{Título}
\author{John Doe}
\date{\today}

\begin{document}
  \maketitle
\end{document}
\end{lstlisting}}
		{\color{Ired} Código}
	\end{column}
	\vline
	\begin{column}{0.5\linewidth}
		{\lmr\centering
		{\LARGE Título \par}
		\vspace{3em}
		{\large\vspace{.75em}
		\begin{tabular}[t]{c}
		John Doe
		\end{tabular}\par}
		\vskip 1.5em
		{\large \today \par}}
		\hfill{\color{Ired} Vista}
	\end{column}
\end{columns}
\end{frame}

\section{Capítulos-Partes, Secciones, etc.}
\begin{frame}{Capítulos-Partes, Secciones, etc.}{}
Al escribir cualquier tipo de documento siempre es bueno llevar una estructura e intentar no romperla a lo largo de este.\pause\\[2em]

En {\lmr\LaTeX} esto se logra mediante los comandos:
\begin{itemize}
	\item \texttt{\textbackslash{}chapter\{<título>\}} para las clases \alert{\bf libro} y \alert{\bf reporte} \\
	\texttt{\textbackslash{}part\{<título>\}} para la clase \alert{\bf articulo}.
	\item \texttt{\textbackslash{}section\{<título>\}}
	\item \texttt{\textbackslash{}subsection\{<título>\}}
	\item Etc.
\end{itemize}
\end{frame}

\begin{frame}{Capítulos-Partes, Secciones, etc.}{}
Siempre que utilizamos estas ordenes {\lmr\LaTeX} las numera de forma automática, teniendo en cuenta que siguen un orden estructural.\pause

\begin{center}
	\begin{tikzpicture}
	\coordinate (zero) at (0,0);
	\node[red,above] (part) at (zero.north) {Parte};
	\node[red,below] (chapter) at (zero.south) {Capítulo};
	
	\node[blue] (sec) at (2,-1.5) {Sección};
	
	\node[cyan] (subsec) at (4,0) {Sub-Sección};
	
	\node[orange] (subsubsec) at (6,-1.5) {Sub-Sub-Sección};
	
	\draw[-latex,shorten <= 1cm,very thick] (zero) -- (sec);
	\draw[-latex,very thick] (sec) -- (subsec);
	\draw[-latex,very thick] (subsec) -- (subsubsec);
	\end{tikzpicture}
\end{center}\pause

Generalmente se suelen utilizar hasta las \alert{\bf sub-secciones}, pero es posible llegar, hasta los \alert{\bf sub-párrafos}.
\end{frame}

\section{El resumen}
\begin{frame}{El Resumen}{}
El resumen es una parte importante de cualquier trabajo escrito.\\[2em]

En {\lmr\LaTeX} existe el entorno \texttt{abstract} para las clases \texttt{article} y \texttt{report}. En este entorno podremos escribir el resumen de nuestro documento, con un margen más pequeño, y con el título que lo identifica.
\end{frame}

\section{Familias, Estilos y Formas}

\section{Tamaño de fuente}

\section{Paquetes {\tt fontenc} e {\tt inputenc}}

\section{{\lmr\LaTeX} en español}

\section{Separación del texto}

\section{Opciones del Documento}
\end{document}
